\subsection{符号の変化が収束を生む場合}

関数が正負に変わると，絶対値を積分した場合と，元の関数を積分した場合とで，収束の結果が異なることがある。

\begin{definition}{絶対収束と条件収束}
広義積分$\int f(x)\dd x$について，$\int|f(x)|\dd x$が収束するとき絶対収束するという。$\int f(x)\dd x$は収束するが，$\int|f(x)|\dd x$は発散するとき，条件収束するという。
\end{definition}

絶対収束なら元の積分も収束する。例えば無限区間では，遠い二点$S<R$の間の積分が
\[
 \left|\int_S^Rf(x)\dd x\right|\leq\int_S^R|f(x)|\dd x
\]
で押さえられる。右辺は$S,R$が十分大きければいくらでも小さくなるので，途中までの積分はコーシーの収束条件を満たす。これは実数の完備性による収束判定であり，「途中までの値が互いにいくらでも近くなるなら，有限の極限をもつ」という事実を使っている。

\begin{example}{振動する積分の条件収束}
$\int_1^{\infty}\frac{\sin x}{x}\dd x$は条件収束することを示せ。
\end{example}

$\frac1x$だけでは無限大で積分可能でない。しかし，$\sin x$を積分すると有界な$-\cos x$になり，$\frac1x$を微分すると$-\frac1{x^2}$になる。この二つを利用するため，部分積分を選ぶ。
\[
\begin{aligned}
 \int_1^R\frac{\sin x}{x}\dd x
 &=\left[-\frac{\cos x}{x}\right]_1^R
   -\int_1^R\frac{\cos x}{x^2}\dd x\\
 &=\cos1-\frac{\cos R}{R}-\int_1^R\frac{\cos x}{x^2}\dd x.
\end{aligned}
\]
境界項$\frac{\cos R}{R}$は$0$へ近づく。また$\frac{|\cos x|}{x^2}\leq\frac1{x^2}$なので，最後の積分は絶対収束する。よって元の積分は収束する。

絶対値を付けた積分は別に調べる。整数$k\geq1$について，$[k\pi,(k+1)\pi]$では$x\leq(k+1)\pi$だから
\[
 \int_{k\pi}^{(k+1)\pi}\frac{|\sin x|}{x}\dd x
 \geq\frac1{(k+1)\pi}\int_{k\pi}^{(k+1)\pi}|\sin x|\dd x
 =\frac2{(k+1)\pi}.
\]
この下限を足すと調和級数の定数倍になり，発散する。したがって絶対値の積分は発散し，元の積分は条件収束である。ここでは積分の値そのものを求めていない。収束を証明することと，値を計算することは別の課題である。

